% !TeX program = lualatex
% Compile twice: lualatex 173.tex
% All references and prose are embedded; no BibTeX database or external inputs.
\documentclass[11pt,a4paper]{article}
\usepackage[margin=24mm,headheight=14pt]{geometry}
\usepackage{fontspec}
\setmainfont{Latin Modern Roman}
\setsansfont{Latin Modern Sans}
\setmonofont{Latin Modern Mono}
\usepackage{amsmath,amssymb,mathtools}
\usepackage{unicode-math}
\setmathfont{Latin Modern Math}
\usepackage{microtype}
\usepackage{enumitem,xurl,needspace}
\usepackage[dvipsnames]{xcolor}
\usepackage{hyperref}
\hypersetup{unicode=true,colorlinks=true,linkcolor=MidnightBlue,urlcolor=MidnightBlue,
 pdftitle={500 Mathematical Problem Statements},pdfauthor={Source-annotated compilation},bookmarksnumbered=true}
\usepackage{bookmark}
\usepackage{tocloft}
\setlength{\cftsecnumwidth}{3em}
\cftsetpnumwidth{2.5em}
\cftsetrmarg{4em}
\usepackage{titlesec}
\titleformat{\section}{\Large\bfseries\raggedright}{\thesection}{0.7em}{}
\usepackage{fancyhdr}
\pagestyle{fancy}\fancyhf{}
\fancyhead[L]{\small\sffamily Mathematical problem statements}
\fancyhead[R]{\small Source edition 22}
\fancyfoot[C]{\thepage}
\setlength{\parindent}{0pt}
\setlength{\parskip}{5pt plus 1pt}
\setlength{\emergencystretch}{3em}
\setcounter{tocdepth}{1}
\setcounter{secnumdepth}{1}
\setlist[itemize]{leftmargin=3em,itemsep=5pt,topsep=4pt,parsep=0pt}
\allowdisplaybreaks
\newcommand{\N}{\mathbb{N}}
\newcommand{\Z}{\mathbb{Z}}
\newcommand{\Q}{\mathbb{Q}}
\newcommand{\R}{\mathbb{R}}
\newcommand{\C}{\mathbb{C}}
\newcommand{\F}{\mathbb{F}}
\newcommand{\A}{\mathbb{A}}
\newcommand{\PP}{\mathbb P}
\newcommand{\E}{\mathbb{E}}
\newcommand{\eps}{\varepsilon}
\newcommand{\dd}{\,\mathrm d}
\newcommand{\sref}[2]{\hyperlink{source:#1:#2}{[S#2]}}
\newcommand{\eref}[2]{\hyperlink{extra:#1:#2}{[E#2]}}
% Turn the next switch off for a shorter reading copy. The quotations preserve
% every exactTarget field, including qualifications omitted from a short summary.
\newif\ifshowcatalogue
\showcataloguetrue
\newcommand{\cataloguescope}[1]{\ifshowcatalogue
 \paragraph{Exact scope in the supplied catalogue.}
 \begin{quote}\small #1\end{quote}\fi}

% Standalone notebook wrapper; original problem section retained verbatim.
\numberwithin{equation}{section}
\hypersetup{pdftitle={Research attempt -- problem 173},
 pdfauthor={Research notebook; source compilation retained}}
\fancyhead[L]{\small\sffamily Research notebook}
\fancyhead[R]{\small\sffamily Problem 173 | 27 September 2026}
\begin{document}
\setcounter{section}{172}
% ================================================================
% problemId: problem.conformal-invariance-of-the-critical-random-cluster-model-with-q-4
% source releaseRank: 173; source releaseStatus: open_with_solved_subcases
% exactTarget SHA-256: bde6c6a02a2b0d5f727bd1f9894109f7d2ad37fa7b26d23000c3f5025a9adf11
\clearpage
\section{\texorpdfstring{Conformal Invariance of the Critical Random-Cluster Model with q <= 4}{Conformal Invariance of the Critical Random-Cluster Model with q <= 4}}\label{173}
\subsection{Definitions and mathematical statement}
On a finite planar graph with boundary condition $\xi$, the FK measure gives a configuration $\omega\subseteq E$ weight proportional to
$p^{|\omega|}(1-p)^{|E|-|\omega|}q^{k_\xi(\omega)}$, where $k_\xi$ counts open clusters after boundary identifications. In a simply connected lattice approximation with marked boundary points, Dobrushin conditions produce an interface between wired primal and dual arcs. The intended statement is convergence of that interface at criticality to chordal $\mathrm{SLE}_{\kappa(q)}$ under mesh refinement. SLE is the Loewner chain driven by $\sqrt\kappa B_t$, and the customary FK relation is $\sqrt q=-2\cos(4\pi/\kappa)$, $4\le\kappa<8$, for $0<q\le4$; see \eref{173}{1}. The catalogue omits the lower range of $q$, lattice, boundary setup, and curve topology. These must be taken from the cited formulation; the formula does not define a probability model for arbitrary negative $q$.
\par\smallskip\textit{Formulation and concept references:} \sref{173}{1}, \eref{173}{1}.
\cataloguescope{Prove that the exploration path of the critical random-cluster (FK) model with q <= 4 converges in distribution to a Schramm-Loewner evolution SLE\_\{kappa(q)\} as the mesh size goes to zero, i.e., the scaling limit is conformally invariant.}
\subsection{Short English statement}
Do critical planar random-cluster interfaces approach the conformally invariant SLE curve predicted by their cluster-weight parameter, with the model and boundary conventions properly fixed?
% BEGIN RESEARCH ADDITION ATTEMPT-173
\subsection{Research attempt}

\paragraph{Current frontier.}
\textbf{Editorial scope note.} The original quotation is retained, including
its omissions. The consulted sources formulate planar FK interfaces with
specific lattice approximations and Dobrushin boundary arcs; a statement
for arbitrary negative $q$ is not a probability assertion. The discussion
below uses the positive-$q$ Dobrushin setup of \sref{173}{1} and
\eref{173}{1}; it does not silently complete every omitted lattice or
topology convention. In particular, FKG-based arguments in the range
$q\ge1$ cannot simply be applied to $0<q<1$.

For $q=2$ the FK-Ising interface and its scaling limit are established;
Kemppainen--Smirnov obtain the full FK-Ising scaling limit in the setting of
\eref{173}{2}. That is not a theorem for all $q\le4$. Nor does triangular
\emph{site} percolation settle every square-lattice FK model at $q=1$.
The exact identification step below is conditional on analytic limits that
are not supplied by a lattice contour identity alone.

\paragraph{Attack route.}
One route proves compactness from crossing estimates and then identifies
every subsequential limit. Another seeks a convergent parafermionic
observable. The first lacks a characterization of the limit; the second
lacks enough regularity and boundary control for general $q$. We combine
them at a concrete point: two continuum martingales, if obtained with the
correct covariance and sufficient integrability, force the entire Loewner
driver. This derives the predicted parameter while making the missing
lattice-to-continuum input explicit.

\paragraph{Attempt: normalization of the candidate observable.}
For $0<q\le4$, choose $0<\sigma\le1$ by
\[
 \sin(\pi\sigma/2)=\sqrt q/2,
 \qquad \kappa=\frac8{1+\sigma}.
\]
Then $4\le\kappa<8$ and
\[
 -2\cos(4\pi/\kappa)
   =-2\cos\bigl(\pi(1+\sigma)/2\bigr)=\sqrt q.
\]
Thus $q=1,2,4$ correspond respectively to
$(\sigma,\kappa)=(1/3,6),(1/2,16/3),(1,4)$, as in the FK parameter
conventions of \eref{173}{1}. This is a normalization check, not a
convergence theorem.

Suppose a subsequential interface limit is generated in the upper half-plane
by continuous chordal Loewner evolution
\[
 \partial_tg_t(z)=\frac2{g_t(z)-U_t},\qquad U_0=0.
\]
For the calculation assume additionally that $U$ is a continuous
semimartingale, with decomposition $U=A+N$, where $A$ is continuous of
finite variation and $N$ is a continuous local martingale. Set $K=[N]$.
These properties are hypotheses here; a continuous driver need not be a
semimartingale merely because it arose as a subsequential limit.
For an interior probe $z$, put
\begin{equation}\label{eq173-observable}
 Z_t(z)=g_t(z)-U_t,
 \qquad M_t(z)=\left(\frac{g_t'(z)}{Z_t(z)}\right)^\sigma.
\end{equation}
Choose analytic branches locally and stop before the probe is swallowed
or the displayed quantities approach zero or infinity. Such stopping
makes every differentiation below legitimate.

\paragraph{Attempt: two martingales force the driver.}
The equations for $Z$ and $g'$ are
\[
 dZ=2Z^{-1}dt-dU,\qquad d\log g'=-2Z^{-2}dt.
\]
It\^o's formula, including the quadratic variation of $Z$, gives
\begin{equation}\label{eq173-ito}
 \frac{dM_t(z)}{M_t(z)}
 =\frac{\sigma}{Z_t(z)}\,dU_t
  +\frac{-4\sigma\,dt+\sigma(1+\sigma)dK_t/2}{Z_t(z)^2}.
\end{equation}
The coefficient $-4\sigma$ includes both the derivative of $g'$ and the
drift in $Z$; omitting either would predict the wrong value of $\kappa$.

Suppose the processes \eqref{eq173-observable} for two distinct probes
$z,z'$ are complex local martingales on a common stopped interval. Uniqueness
of the semimartingale decomposition says their finite-variation parts
vanish. Multiplication by the bounded nonzero factors gives equalities of
complex measures
\begin{equation}\label{eq173-measures}
 Z_t(z)\,dA_t+\frac{1+\sigma}{2}\,dK_t-4dt=0,
 \qquad
 Z_t(z')\,dA_t+\frac{1+\sigma}{2}\,dK_t-4dt=0.
\end{equation}
Before either probe is swallowed, conformality implies
$g_t(z)\ne g_t(z')$. Subtracting the equations yields $dA=0$; locally
one may divide by their nonzero continuous difference. Substitution yields
$dK=8dt/(1+\sigma)=\kappa dt$. L\'evy's characterization now gives
$U=\sqrt\kappa B$ on that interval. An exhaustion by stopped intervals,
with suitable probes covering every finite time, gives the chordal
$\mathrm{SLE}_\kappa$ driver.

This is a proved stochastic identification lemma under its stated
hypotheses. It does not assert that the particular FK observable has
already converged to \eqref{eq173-observable}. Two fixed probes can be
swallowed; the exhaustion or a sufficiently rich family of probes is
necessary for a conclusion about all times.

\paragraph{Attempt: the lattice-to-limit inputs that would close the route.}
Here is a sufficient package for a fixed admissible lattice, domain, and
parameter $q$. First obtain tightness of the interfaces in the chosen curve
topology and enough control to pass to Loewner chains parametrized by
half-plane capacity. Next show that the normalized lattice observable,
with its prescribed boundary values, converges to
\eqref{eq173-observable} along every subsequence, jointly with the explored
hulls. Finally justify passage of its conditional martingale identities
and establish a semimartingale driver, or derive that property by an
independent characterization theorem. The preceding lemma then identifies
every subsequential law. Tightness and uniqueness of all limit laws imply
convergence of the full sequence.

A usable martingale-passage argument must be quantitative. On each stopped
interval, joint distributional convergence together with uniform
integrability allows passage of
$\E[(M_t-M_s)F]=0$ for bounded continuous functions $F$ of the past.
It then extends to the limiting past sigma-field by approximation. Without
uniform integrability, convergence in distribution of martingales need
not preserve their conditional means. Stopping a proposed continuum
observable does not automatically bound its lattice precursors, so a
separate lattice estimate is still required.

The discrete parafermionic contour relation described in \eref{173}{1}
is a motivation for this package, not a substitute for it. To identify a
holomorphic boundary-value solution, one needs precompactness, the correct
boundary behavior, and uniqueness in a suitable analytic class. For general
$q$, a local linear contour identity alone has not supplied those inputs
in this attempt. At $q=4$ the stochastic calculation is nonsingular, but
that observation does not remove possible logarithmic-scale difficulties
in the lattice estimates.

\paragraph{Stress test.}
No use was made of independent increments before deriving Brownian motion.
The semimartingale assumption, the uniform-integrability requirement, and
the availability of probes after successive swallowing times have not been
hidden in the word ``martingale.'' The computations are stopped away from
singularities, so fractional powers and It\^o derivatives are well defined.

The limit lemma works for each positive $q$ with the displayed parameter,
but it does not manufacture a critical FK model for omitted or inadmissible
conventions. In particular, results proved only for FK-Ising
\eref{173}{2} are not invoked for other cluster weights. The extra geometric
control needed for trace convergence is not supplied solely by convergence
of driving functions, and the topology in the original source must remain
part of any completed statement.

\paragraph{Outcome.}
\textbf{Outcome: Conditional reduction.}

The attempt proves a two-probe continuum identification lemma and derives
the exact FK--SLE parameter relation from its drift cancellation. The
missing input is a lattice convergence and tightness package strong enough
to realize those martingales, with the specified boundary and topology
conventions. The argument is not a new general conformal-invariance theorem;
its standard Loewner and stochastic ingredients are used to make the gap
explicit rather than to assume it away.
% END RESEARCH ADDITION ATTEMPT-173
\subsection{Sources}
\begin{itemize}\raggedright
\item[\textnormal{[S1]}]\hypertarget{source:173:1}{} H. Duminil-Copin, 'Lectures on the Ising and Potts models on the hypercubic lattice', arXiv:1707.00520 (2017), in Random Graphs, Phase Transitions and the Gaussian Free Field, Springer PIMS 123 (2020).
\url{https://ar5iv.labs.arxiv.org/html/1707.00520}.
{\small\textit{Catalogue formulation source.}}
\item[\textnormal{[E1]}]\hypertarget{extra:173:1}{} Hugo Duminil-Copin and Stanislav Smirnov, Conformal invariance of lattice models, arXiv:1109.1549v4, FK interface and SLE discussion.
\url{https://arxiv.org/html/1109.1549v4}.
{\small\textit{Added primary source: the random-cluster interface setup and parameter conventions. Consulted 24 September 2026.}}
% BEGIN RESEARCH ADDITION SOURCES-173
\item[\textnormal{[E2]}]\hypertarget{extra:173:2}{} Antti Kemppainen and
Stanislav Smirnov, \emph{Conformal invariance in random cluster models.
II. Full scaling limit as a branching SLE}, arXiv:1609.08527, 2016,
abstract and FK-Ising scaling-limit theorem.
\url{https://arxiv.org/abs/1609.08527}.
{\small\textit{Consulted 27 September 2026. The cited result concerns
$q=2$, not arbitrary $q\le4$.}}
% END RESEARCH ADDITION SOURCES-173
\end{itemize}

\end{document}
